\ifdefined\gkver\else\def\gkver{student}\fi
\documentclass[\gkver]{gaokaozhenti}
\begin{document}

\chapter{2024年江苏}

\begin{enumerate}
\item
%% number: 1
%% paperName: 2024年江苏高考第 1 题 4 分
%% typeId: 1
%% type: 单选题
%% chapter: 电场
%% point: 电场中的图象
%% method: 无
%% score: 4
%% degree: 950
%% duplicateId: 0
%% body:
在静电场中有 a、b 两点，试探电荷在两点的静电力 $F$ 与电荷量 $q$ 满足如图所示的关系，请问 a、b 两点的场强大小之比 $E_{a}:E_{b}$ 为\xzanswer{D}

\onepicture[w=3.40cm]{figs/01.png}

\fourchoices[answer=D]
{$1:1$}
{$2:1$}
{$3:1$}
{$4:1$}

%% answer: D
%% memo:
\memoanswer{
根据 $E = \frac{F}{q}$ 可知 $F - q$ 图像斜率表示电场强度，由图可知 a、b 两图线的斜率之比为 $4:1$，即 $E_{a}:E_{b} = 4:1$。故选 D。
}

\item
%% number: 2
%% paperName: 2024年江苏高考第 2 题 4 分
%% typeId: 1
%% type: 单选题
%% chapter: 光学
%% point: 光的偏振
%% method: 无
%% score: 4
%% degree: 950
%% duplicateId: 0
%% body:
用立体影院的特殊眼镜去观看手机液晶屏幕，左镜片明亮，右镜片暗，现在将手机屏幕旋转 90 度，会观察到\xzanswer{D}

\fourchoices[answer=D]
{两镜片都变亮}
{两镜片都变暗}
{两镜片没有任何变化}
{左镜片变暗，右镜片变亮}

%% answer: D
%% memo:
\memoanswer{
立体影院的特殊眼镜是利用了光的偏振，其镜片为偏振片。用立体影院的特殊眼镜去观看手机液晶屏幕，左镜片明亮，右镜片暗，根据 $I = I_{0}\cos^{2}\theta$ 可知将手机屏幕旋转 90 度后左镜片变暗，右镜片变亮。

故选 D。
}

\item
%% number: 3
%% paperName: 2024年江苏高考第 3 题 4 分
%% typeId: 1
%% type: 单选题
%% chapter: 原子和原子核
%% point: 人工核反应
%% method: 无
%% score: 4
%% degree: 950
%% duplicateId: 0
%% body:
用粒子轰击氮核从原子核中打出了质子，该实验的核反应方程式是 $\ce{X + ^{14}_{7}N -> ^{1}_{1}H + ^{17}_{8}O}$，粒子 X 为\xzanswer{D}

\fourchoices[answer=D]
{正电子 $\ce{^{0}_{1}e}$}
{中子 $\ce{^{1}_{0}n}$}
{氘核 $\ce{^{2}_{1}H}$}
{氦核 $\ce{^{4}_{2}He}$}

%% answer: D
%% memo:
\memoanswer{
根据质量数守恒可知 X 的质量数为 $m = 17 + 1 - 14 = 4$

根据电荷守恒可知 X 的电荷数为 $n = 8 + 1 - 7 = 2$

可知 X 为 $\alpha$ 粒子 $\ce{^{4}_{2}He}$。

故选 D。
}

\item
%% number: 4
%% paperName: 2024年江苏高考第 4 题 4 分
%% typeId: 1
%% type: 单选题
%% chapter: 曲线运动
%% point: 平抛运动
%% method: 无
%% score: 4
%% degree: 750
%% duplicateId: 0
%% body:
喷泉 a、b 形成如图所示的形状，不计空气阻力，则喷泉 a、b 的\xzanswer{A}

\onepicture[w=4.20cm]{figs/04.png}

\fourchoices[answer=A]
{加速度相同}
{初速度相同}
{最高点的速度相同}
{在空中的时间相同}

%% answer: A
%% memo:
\memoanswer{
A．不计空气阻力，在喷泉喷出的水在空中只受重力，加速度均为重力加速度，故 A 正确；

D．设喷泉喷出的水竖直方向的分速度为 $v_{y}$，水平方向速度为 $v_{x}$，竖直方向，根据对称性可知在空中运动的时间 $t = 2\sqrt{\frac{2h}{g}}$ 可知 $t_{b} > t_{a}$。D 错误；

BC．最高点的速度等于水平方向的分速度 $v_{x}=\frac{x}{t}$，由于水平方向的位移大小关系未知，无法判断最高点的速度大小关系，根据速度的合成可知无法判断初速度的大小，BC 错误；故选 A。
}

\item
%% number: 5
%% paperName: 2024年江苏高考第 5 题 4 分
%% typeId: 1
%% type: 单选题
%% chapter: 原子和原子核
%% point: 玻尔模型
%% method: 无
%% score: 4
%% degree: 850
%% duplicateId: 0
%% body:
在原子跃迁中，辐射如图所示的 4 种光子，其中只有一种光子可使某金属发生光电效应，是哪一种\xzanswer{C}

\onepicture[w=4.00cm]{figs/05.png}

\fourchoices[answer=C]
{$\lambda_{1}$}
{$\lambda_{2}$}
{$\lambda_{3}$}
{$\lambda_{4}$}

%% answer: C
%% memo:
\memoanswer{
根据光电方程可知当只有一种光子可使某金属发生光电效应该光子对应的能量最大，根据图中能级图可知跃迁时对应波长为 $\lambda_{3}$ 的光子能量最大。

故选 C。
}

\item
%% number: 6
%% paperName: 2024年江苏高考第 6 题 4 分
%% typeId: 1
%% type: 单选题
%% chapter: 光学
%% point: 光的折射
%% method: 无
%% score: 4
%% degree: 750
%% duplicateId: 0
%% body:
现有一光线以相同的入射角 $\theta$，打在不同浓度 NaCl 的两杯溶液中，折射光线如图所示（$\beta_{1} < \beta_{2}$），已知折射率随浓度增大而变大。则\xzanswer{A}

\onepicture[w=4.60cm]{figs/06.png}

\fourchoices[answer=A]
{甲折射率大}
{甲浓度小}
{甲速度大}
{甲临界角大}

%% answer: A
%% memo:
\memoanswer{
入射角相同，由于 $\beta_{1} < \beta_{2}$，根据折射定律可知甲溶液折射率大于乙溶液，故甲浓度大；根据 $v=\frac{c}{n}$，可知光线在甲中的传播速度较小，由 $\sin C=\frac{1}{n}$ 可知折射率越大临界角越小，故甲临界角小。

故选 A。
}

\item
%% number: 7
%% paperName: 2024年江苏高考第 7 题 4 分
%% typeId: 1
%% type: 单选题
%% chapter: 机械波
%% point: 波长频率波速
%% method: 无
%% score: 4
%% degree: 650
%% duplicateId: 0
%% body:
如图所示，水面上有 O、A、B 三点共线，OA = 2AB，$t = 0$ 时刻在 O 点的水面给一个扰动，$t_{1}$ 时刻 A 开始振动，则 B 振动的时刻为\xzanswer{B}

\onepicture[w=5.40cm]{figs/07.png}

\fourchoices[answer=B]
{$t_{1}$}
{$\frac{3}{2}t_{1}$}
{$2t_{1}$}
{$\frac{5}{2}t_{1}$}

%% answer: B
%% memo:
\memoanswer{
机械波的波速 $v$ 不变，设 $OA = 2AB = 2L$，故可得 $t_{1} = \frac{2L}{v}$

可得 $t_{AB}=\frac{L}{v}=\frac{1}{2}t_{1}$

故可得 B 振动的时刻为 $t = t_{1} + t_{AB} = \frac{3}{2}t_{1}$

故选 B。
}

\item
%% number: 8
%% paperName: 2024年江苏高考第 8 题 4 分
%% typeId: 1
%% type: 单选题
%% chapter: 曲线运动
%% point: 圆周运动的应用
%% method: 无
%% score: 4
%% degree: 650
%% duplicateId: 0
%% body:
陶瓷是以粘土为主要原料以及各种天然矿物经过粉碎混炼、成型和煅烧制得的材料以及各种制品。如图所示是生产陶磁的简化工作台，当陶磁匀速转动时，台面上掉有陶屑，陶屑与桌面间的动摩擦因数处处相同（台面够大），则\xzanswer{D}

\onepicture[w=3.60cm]{figs/08.png}

\fourchoices[answer=D]
{离轴 $OO'$ 越远的陶屑质量越大}
{离轴 $OO'$ 越近的陶屑质量越小}
{只有平台边缘有陶屑}
{离轴最远的陶屑距离不会超过某一值}

%% answer: D
%% memo:
\memoanswer{
ABC．与台面相对静止的陶屑做匀速圆周运动，静摩擦力提供向心力，当静摩擦力为最大静摩擦力时，根据牛顿第二定律可得 $\mu mg = m\omega^{2}r$

解得 $r = \frac{\mu g}{\omega^{2}}$

因与台面相对静止的这些陶屑的角速度相同，由此可知能与台面相对静止的陶屑离轴 $OO'$ 的距离与陶屑质量无关，只要在台面上不发生相对滑动的位置都有陶屑。故 ABC 错误；

D．离轴最远的陶屑其受到的静摩擦力为最大静摩擦力，由前述分析可知最大的运动半径为 $R = \frac{\mu g}{\omega^{2}}$

$\mu$ 与 $\omega$ 均一定，故 $R$ 为定值，即离轴最远的陶屑距离不超过某一值 $R$。故 D 正确。

故选 D。
}

\item
%% number: 9
%% paperName: 2024年江苏高考第 9 题 4 分
%% typeId: 1
%% type: 单选题
%% chapter: 运动和力
%% point: 动量和能量
%% method: 无
%% score: 4
%% degree: 550
%% duplicateId: 0
%% body:
在水平面上有一个 U 形滑板 A，A 的上表面有一个静止的物体 B，左侧用轻弹簧连接在物体 B 的左侧，右侧用一根细绳连接在物体 B 的右侧，开始时弹簧处于拉伸状态，各表面均光滑，剪断细绳后，则\xzanswer{A}

\onepicture[w=4.40cm]{figs/09.png}

\fourchoices[answer=A]
{弹簧原长时 B 动量最大}
{压缩最短时 A 动能最大}
{系统动量变大}
{系统机械能变大}

%% answer: A
%% memo:
\memoanswer{
对整个系统分析可知合外力为 0，A 和 B 组成的系统动量守恒，得

$m_{A}v_{A}=m_{B}v_{B}$

设弹簧的初始弹性势能为 $E_{p}$，整个系统只有弹簧弹力做功，机械能守恒，当弹簧原长时得

$E_{\mathrm{p}}=\frac{1}{2}m_{\mathrm{A}}v_{\mathrm{A}}^{2}+\frac{1}{2}m_{\mathrm{B}}v_{\mathrm{B}}^{2}$

联立得

$E_{\mathrm{p}}=\frac{1}{2}\left(\frac{m_{\mathrm{B}}^{2}}{m_{\mathrm{A}}}+m_{\mathrm{B}}\right)v_{\mathrm{B}}^{2}$

故可知弹簧原长时物体速度最大，此时动量最大，动能最大。对于系统来说动量一直为零，系统机械能不变。

故选 A。
}

\item
%% number: 10
%% paperName: 2024年江苏高考第 10 题 0 分
%% typeId: 1
%% type: 单选题
%% chapter: 电磁感应
%% point: 楞次定律
%% method: 无
%% score: 0
%% degree: 550
%% duplicateId: 0
%% body:
如图所示，在绝缘的水平面上，有闭合的两个线圈 a、b，线圈 a 处在匀强磁场中，现将线圈 a 从磁场中匀速拉出，线圈 a、b 中产生的感应电流方向分别是\xzanswer{A}

\onepicture[w=4.20cm]{figs/10.png}

\fourchoices[answer=A]
{顺时针，顺时针}
{顺时针，逆时针}
{逆时针，顺时针}
{逆时针，逆时针}

%% answer: A
%% memo:
\memoanswer{
线圈 a 从磁场中匀速拉出的过程中穿过 a 线圈的磁通量在减小，则根据楞次定律可知 a 线圈的电流为顺时针，由于线圈 a 从磁场中匀速拉出则 a 中产生的电流为恒定电流，则线圈 a 靠近线圈 b 的过程中线圈 b 的磁通量在向外增大，同理可得线圈 b 产生的电流为顺时针。故选 A。
}

\item
%% number: 11
%% paperName: 2024年江苏高考第 11 题 4 分
%% typeId: 1
%% type: 单选题
%% chapter: 曲线运动
%% point: 圆周运动的应用
%% method: 无
%% score: 4
%% degree: 450
%% duplicateId: 0
%% body:
如图所示，细绳穿过竖直的管子拴住一个小球，让小球在 A 高度处作水平面内的匀速圆周运动，现用力将细绳缓慢下拉，使小球在 B 高度处作水平面内的匀速圆周运动，不计一切摩擦，则\xzanswer{BD}

\onepicture[w=2.60cm]{figs/11.png}

\fourchoices[answer=BD]
{周期 $T_{A} < T_{B}$}
{角速度 $\omega_{A} < \omega_{B}$}
{线速度 $v_{A} > v_{B}$}
{向心加速度 $a_{A} < a_{B}$}

%% answer: BD
%% memo:
\memoanswer{
A．设绳子与竖直方向夹角为 $\theta$，绳子长度为 $l$，小球所在平面距离顶点的竖直高度为 $h$，对小球分析有

$mg \tan \theta = m \frac{4\pi^{2}}{T^{2}} l \sin \theta$

整理有

$T=2\pi\sqrt{\frac{l\cos\theta}{g}}=2\pi\sqrt{\frac{h}{g}}$

由题意可知，高度 $h$ 变小，所以周期变小，即 $T_{A} > T_{B}$，故 A 项错误；

B．对小球分析有 $mg \tan \theta = m \omega^{2} l \sin \theta$

整理有

$\omega=\sqrt{\frac{g}{l\cos\theta}}=\sqrt{\frac{g}{h}}$

由题意可知，高度 $h$ 变小，所以角速度大小变大，即 $\omega_{A} < \omega_{B}$，故 B 项正确；

C．对小球分析有 $mg \tan \theta = m \frac{v^{2}}{l \sin \theta}$

整理有

$v = \sin \theta \sqrt{\frac{gl}{\cos \theta}}$

由题意可知，绳子长度 $l$ 变小，角度 $\theta$ 变大，所以线速度大小变化无法判断，故 C 项错误；

D．对小球分析有 $mg \tan \theta = ma$

整理有 $a = g \tan \theta$

由题意可知，其角度 $\theta$ 变大，所以向心加速度大小变大，即 $a_{A} < a_{B}$，故 D 项正确。故选 BD。
}

\item
%% number: 12
%% paperName: 2024年江苏高考第 12 题 15 分
%% typeId: 4
%% type: 实验题
%% chapter: 电路
%% point: 测量电阻率
%% method: 无
%% score: 15
%% degree: 550
%% duplicateId: 0
%% body:
有一块长方体霍尔元件，长、宽、高分别为 $a$、$b$、$c$，如图~\ref{2024江苏12a}所示。为了测量该霍尔元件的电阻率，进行了如下操作。
\begin{enumerate}
	\item
	用多用电表测量电阻，沿 ab 方向测得的电阻为 $10 \UO$，沿 bc 方向的电阻如图~\ref{2024江苏12b}所示，由图读出沿 bc 方向的电阻为\tkanswer[1.6cm]{300}$\UO$；
	\item
	某同学根据如图~\ref{2024江苏12c}所示的电路图连接实物图丁，请判断连接错误的区域是\tkanswer[1.2cm]{A}；
	\item
	测量 bc 方向的电阻时，手头有两个滑动变阻器，应选择\tkanswer[0.8cm]{B}；

	A．滑动变阻器：最大阻值为 $5 \UO$，允许通过的最大电流为 $1.0 \UA$

	B．滑动变阻器：最大阻值为 $500 \UO$，允许通过的最大电流为 $0.5 \UA$
	\item
	接通开关前，滑动变阻器滑片应放在\tkanswer[1.6cm]{右端}；
	\item
	测量小电阻时，用微安表（量程 $0 \sim 100 \UuA$，内阻约为 $4 \UO$），测得电阻率为 $1.15 \UOm$，测量大电阻时，用电流表（量程 $0 \sim 100 \UmA$，内阻约为 $1 \UO$），测得电阻率为 $1.32 \UOm$，小明说，沿 ab 方向的电阻小，所测量的误差小，请判断是否正确？简述理由\tkanswer[4.4cm]{错误；沿 ab 方向的电阻约为 $10 \UO$，若用内阻约为 $4 \UO$ 的微安表测量，其误差要比内阻约为 $1 \UO$ 的电流表测量较大的 bc 间的约 $300 \UO$ 电阻误差肯定较大，则小明的说法错误}。
\end{enumerate}

\fourpicture[labelA=2024江苏12a, labelB=2024江苏12b, labelC=2024江苏12c, labelD=2024江苏12d, numA=甲, numB=乙, numC=丙, numD=丁, wA=3.20cm, wB=3.00cm, wC=3.00cm, wD=3.60cm, align=right]
{figs/12a.png}
{figs/12b.png}
{figs/12c.png}
{figs/12d.png}
%% answer:
\jdanswer{
\begin{enumerate}
	\item
	$300$

	\item
	A

	\item
	B

	\item
	右端

	\item
	错误；沿 ab 方向的电阻约为 $10 \UO$，若用内阻约为 $4 \UO$ 的微安表测量，其误差要比内阻约为 $1 \UO$ 电流表测量较大的 bc 间的约 $300 \UO$ 电阻误差肯定较大，则小明的说法错误。

\end{enumerate}
}
%% memo:
\memoanswer{
（1）沿 bc 方向的电阻为 $3 \times 100 \UO = 300 \UO$；

（2）连接错误的区域是 A 区域，电压表的左边接线柱应该连接到电阻的左端；

（3）因 bc 方向的电阻约为 $300 \UO$，则测量 bc 方向的电阻时，应该选择与待测电阻阻值相当的滑动变阻器 B 即可；

（4）接通开关前，滑动变阻器滑片应放在阻值最大的最右端位置；

（5）沿 ab 方向的电阻约为 $10 \UO$，若用内阻约为 $4 \UO$ 的微安表测量，其误差要比内阻约为 $1 \UO$ 电流表测量较大的 bc 间的约 $300 \UO$ 电阻误差肯定较大，则小明的说法错误。
}

\item
%% number: 13
%% paperName: 2024年江苏高考第 13 题 6 分
%% typeId: 6
%% type: 计算题
%% chapter: 气体性质
%% point: 查理定律
%% method: 无
%% score: 6
%% degree: 750
%% duplicateId: 0
%% body:
某科研实验站有一个密闭容器，容器内有温度为 $300 \UK$，压强为 $10^{5} \UPa$ 的气体，容器内有一个面积为 $0.06 \Umq$ 的观测台，现将这个容器移动到月球，容器内的温度变成 $240 \UK$，整个过程可认为气体的体积不变，月球表面为真空状态。求：
\begin{enumerate}
	\item
	气体现在的压强；
	\item
	观测台对气体的压力。
\end{enumerate}
%% answer:
\jdanswer{
\begin{enumerate}
	\item
	$8\times10^{4} \UPa$

	\item
	$4.8\times10^{3} \UN$

\end{enumerate}
}
%% memo:
\memoanswer{
（1）由题知，整个过程可认为气体的体积不变，则有

$\frac{p_{1}}{T_{1}} = \frac{p_{2}}{T_{2}}$

解得 $p_{2} = 8\times10^{4} \UPa$

（2）根据压强的定义，观测台对气体的压力

$F = p_{2}S = 4.8\times10^{3} \UN$
}

\item
%% number: 14
%% paperName: 2024年江苏高考第 14 题 8 分
%% typeId: 6
%% type: 计算题
%% chapter: 运动和力
%% point: 动量守恒定律
%% method: 无
%% score: 8
%% degree: 650
%% duplicateId: 0
%% body:
嫦娥六号在轨速度为 $v_{0}$，着陆器对应的组合体 A 与轨道器对应的组合体 B 分离时间为 $\Delta t$，分离后 B 的速度为 $v$，且与 $v_{0}$ 同向，A、B 的质量分别为 $m$、$M$。求：
\begin{enumerate}
	\item
	分离后 A 的速度 $v_{1}$ 大小；
	\item
	分离时 A 对 B 的推力大小。
\end{enumerate}
%% answer:
\jdanswer{
\begin{enumerate}
	\item
	$\frac{(m + M)v_{0} - Mv}{m}$

	\item
	$\frac{M(v - v_{0})}{\Delta t}$

\end{enumerate}
}
%% memo:
\memoanswer{
（1）组合体分离前后动量守恒，取 $v_{0}$ 的方向为正方向，有

$(m+M)v_{0} = Mv + mv_{1}$

解得 $v_{1} = \frac{(m + M)v_{0} - Mv}{m}$

方向与 $v_{0}$ 相同；

（2）以 B 为研究对象，对 B 列动量定理有

$F\Delta t = Mv - Mv_{0}$

解得 $F = \frac{M(v - v_{0})}{\Delta t}$
}

\item
%% number: 15
%% paperName: 2024年江苏高考第 15 题 12 分
%% typeId: 6
%% type: 计算题
%% chapter: 功和能
%% point: 功能原理
%% method: 无
%% score: 12
%% degree: 550
%% duplicateId: 0
%% body:
如图所示，粗糙斜面的动摩擦因数为 $\mu$，倾角为 $\theta$，斜面长为 $L$。一个质量为 $m$ 的物块，在电动机作用下，从 A 点由静止加速至 B 点时达到最大速度 $v$，之后作匀速运动至 C 点，关闭电动机，从 C 点又恰好到达最高点 D。求：
\begin{enumerate}
	\item
	CD 段长 $x$；
	\item
	BC 段电动机的输出功率 $P$；
	\item
	全过程物块增加的机械能 $E_{1}$ 和电动机消耗的总电能 $E_{2}$ 的比值。
\end{enumerate}
\onepicture[w=5.20cm, align=right]{figs/15.png}
%% answer:
\jdanswer{
\begin{enumerate}
	\item
	$\frac{v^{2}}{2g(\sin \theta + \mu \cos \theta)}$

	\item
	$mgv(\sin\theta + \mu\cos\theta)$

	\item
	$\frac{\sin \theta}{\sin \theta + \mu \cos \theta}$

\end{enumerate}
}
%% memo:
\memoanswer{
（1）物块在 CD 段运动过程中，由牛顿第二定律得

$mg\sin\theta + \mu mg\cos\theta = ma$

由运动学公式 $0 - v^{2} = -2ax$

联立解得 $x = \frac{v^{2}}{2g(\sin\theta + \mu\cos\theta)}$

（2）物块在 BC 段匀速运动，得电动机的牵引力为

$F = mg\sin\theta + \mu mg\cos\theta$

由 $P = Fv$ 得

$P = mgv(\sin\theta + \mu\cos\theta)$

（3）全过程物块增加的机械能为 $E_{1} = mgL\sin\theta$

整个过程由能量守恒得电动机消耗的总电能转化为物块增加的机械能和摩擦产生的内能，故可知 $E_{2} = E_{1} + \mu mg\cos\theta\cdot L$

故可得 $\frac{E_{1}}{E_{2}} = \frac{mgL\sin\theta}{mgL\sin\theta + \mu mgL\cos\theta} = \frac{\sin\theta}{\sin\theta + \mu\cos\theta}$
}

\item
%% number: 16
%% paperName: 2024年江苏高考第 16 题 15 分
%% typeId: 6
%% type: 计算题
%% chapter: 磁场
%% point: 带电粒子在组合场中的运动
%% method: 无
%% score: 15
%% degree: 350
%% duplicateId: 0
%% body:
如图所示，两个半圆区域 abcd、a$'$b$'$c$'$d$'$ 中有垂直纸面向里的匀强磁场，半径分别为 $R_{1}$、$R_{2}$。ab 与 a$'$b$'$ 间有一个匀强电场，电势差为 $U$，cd 与 c$'$d$'$ 间有一个插入体，电子每次经过插入体速度减为原来的 $k$ 倍。现有一个质量为 $m$、电量为 $e$ 的电子，从 cd 面射入插入体，经过磁场、电场后再次到达 cd 面，速度增加，多次循环运动后，电子的速度大小达到一个稳定值，忽略相对论效应，忽略经过插入体的时间。求：
\begin{enumerate}
	\item
	电子进入插入体前后在磁场中的半径 $r_{1}$、$r_{2}$ 之比；
	\item
	电子多次循环后到达 cd 的稳定速度 $v$；
	\item
	若电子到达 cd 中点 P 时速度稳定，并最终到达边界 d，求电子从 P 到 d 的时间 $t$。
\end{enumerate}
\onepicture[w=4.60cm, align=right]{figs/16.png}
%% answer:
\jdanswer{
\begin{enumerate}
	\item
	$\frac{1}{k}$

	\item
	$v = \sqrt{\frac{2eU}{(1 - k^{2})m}}$

	\item
	$t = \frac{\pi(R_{2} - R_{1})}{2}\sqrt{\frac{(1 + k)m}{2(1 - k)eU}}$

\end{enumerate}
}
%% memo:
\memoanswer{
（1）根据洛伦兹力提供向心力，可得 $Bev = m \frac{v^{2}}{r}$

解得 $r = \frac{mv}{Be}$

设电子进入插入体前的速度为 $v_{1}$，穿出插入体后的速度为 $v_{2}$，可得

$\frac{r_{1}}{r_{2}} = \frac{v_{1}}{v_{2}} = \frac{v_{1}}{kv_{1}} = \frac{1}{k}$

（2）根据题意可得 $\frac{1}{2}mv^{2} - \frac{1}{2}m(kv)^{2} = eU$

解得 $v = \sqrt{\frac{2eU}{(1 - k^{2})m}}$

（3）电子到达 cd 中点 P 时速度稳定，电子在右侧区域运动时速度为 $v$，做匀速圆周运动的半径 $r_{\text{右}} = \frac{mv}{Be}$

电子在左侧区域运动时速度为 $kv$，做匀速圆周运动的半径 $r_{\text{左}} = \frac{kmv}{Be} = kr_{\text{右}}$

设电子从 P 点转过 $n$ 个圆周到达 d 点，根据题意可得 $n(2r_{\text{右}} - 2r_{\text{左}}) = \frac{1}{2}(R_{2} - R_{1})$

可得 $n = \frac{R_{2} - R_{1}}{4(r_{\text{右}} - r_{\text{左}})}$

运动时间 $t = nT$，其中 $T = \frac{2\pi m}{Be}$

整理可得 $t = \frac{R_{2} - R_{1}}{4(r_{\text{右}} - r_{\text{左}})} \cdot \frac{2\pi m}{Be} = \frac{\pi(R_{2} - R_{1})}{2v(1 - k)} = \pi(R_{2} - R_{1})\sqrt{\frac{(1 + k)m}{2(1 - k)eU}}$
}

\end{enumerate}

\end{document}
